% 250,120 Spectral Anomalies from DESI DR1: A 173-Column Enhanced Catalog
% from the Largest Autoencoder Search at Survey Scale
% Houston Golden — 2026
% Target journal: ApJS
%
% Compile: pdflatex paper3_anomaly_catalog && bibtex paper3_anomaly_catalog && pdflatex paper3_anomaly_catalog && pdflatex paper3_anomaly_catalog

\documentclass[aps,prd,twocolumn,superscriptaddress,showpacs,preprintnumbers,nofootinbib,longbibliography,floatfix]{revtex4-2}

\usepackage{amsmath,amssymb,amsfonts}
\usepackage{graphicx}
\usepackage{bm}
\usepackage{hyperref}
\usepackage{xcolor}
\usepackage{booktabs}
\usepackage{multirow}
\usepackage{dcolumn}
\usepackage{enumitem}
\usepackage{mathtools}
\usepackage[utf8]{inputenc}
\usepackage{float}
\usepackage{longtable}

\hypersetup{
  colorlinks=true,
  linkcolor=blue!60!black,
  citecolor=blue!60!black,
  urlcolor=blue!60!black
}

% Custom commands
\newcommand{\bigae}{\textsc{BigAE}}
\newcommand{\desi}{\textsc{DESI}}
\newcommand{\redrock}{\textsc{Redrock}}
\newcommand{\simbad}{\textsc{SIMBAD}}
\newcommand{\hdbscan}{\textsc{HDBSCAN}}
\newcommand{\umap}{\textsc{UMAP}}
\newcommand{\healpix}{\textsc{HEALPix}}
\newcommand{\etal}{\textit{et~al.}}
\newcommand{\ie}{i.e.}
\newcommand{\eg}{e.g.}
\newcommand{\cf}{cf.}
\newcommand{\snr}{\mathrm{SNR}}
\newcommand{\zwarn}{\texttt{ZWARN}}
\newcommand{\paperVersion}{v0.1.0}
\newcommand{\paperTimestamp}{2026-03-28}

\begin{document}

\preprint{}

\title{250,120 Spectral Anomalies from DESI DR1:\\
A 173-Column Enhanced Catalog from the Largest\\
Autoencoder Search at Survey Scale}

\author{Houston Golden}
\email{houston@hubify.com}
\affiliation{Independent Researcher, Los Angeles, California, USA}

\date{March 28, 2026 --- \paperVersion}

% ============================================================
% ABSTRACT
% ============================================================
\begin{abstract}
We present an enhanced catalog of 22,504,897 \desi{} Data Release~1 spectra scored by an unsupervised autoencoder (\bigae{}, $496 \to 128 \to 496$, ${\sim}660$K parameters).
Each spectrum is profiled with 173 columns including 128-dimensional latent vectors, per-band anomaly scores, pipeline classifications, photometry, and derived features.
We identify 83 high-confidence anomalies (anomaly score $> 3\sigma$, $\snr > 2$, $\zwarn = 0$), dominated by 69 quasi-stellar objects (QSOs) at mean redshift $\bar{z} = 5.5$, including 12 reionization-era QSOs at $z > 6$ exhibiting Gunn--Peterson troughs and unusual emission features.
\umap{}$+$\hdbscan{} clustering of the 128-dimensional latent space reveals two distinct anomaly populations: a dominant low-$\snr$ cluster and 2,575 ``red-anomaly'' objects presenting the paradox of high pipeline confidence yet high autoencoder reconstruction error.
The anomaly rate is uniform (${\sim}1\%$) across the \desi{} footprint (Spearman $r = 0.03$ with survey depth).
We honestly disclose that the bulk anomaly count (250,120 at ${>}5\sigma$) is dominated by low-$\snr$ spectra; the scientifically actionable catalog contains 83--115 objects.
The enhanced catalog with latent vectors enables downstream applications including spectral taxonomy, photometric redshift estimation, and cross-survey correlation.
The catalog and trained model are publicly available on HuggingFace.
\end{abstract}

\maketitle

% ============================================================
% 1. INTRODUCTION
% ============================================================
\section{Introduction}
\label{sec:intro}

The Dark Energy Spectroscopic Instrument (\desi{}) is conducting the largest spectroscopic survey in history, with a design goal of ${\sim}40$ million galaxy and quasar redshifts over 14,000~deg$^2$ \cite{DESI2024}.
Data Release~1 (DR1) provides ${\sim}22.5$ million spectra obtained during the first 14 months of main survey operations, spanning diverse target classes including luminous red galaxies (LRGs), emission-line galaxies (ELGs), bright galaxies (BGS), QSOs, stars, and Milky Way Survey (MWS) targets \cite{DESI2025DR2}.
The sheer volume of \desi{} spectra guarantees that rare, unusual, and previously uncataloged objects are present in the data --- objects that the standard \redrock{} pipeline \cite{Guy2023} may classify correctly by redshift but whose spectral morphology departs significantly from the template basis.

Anomaly detection in large spectroscopic surveys has emerged as a powerful technique for discovering unusual objects that do not fit standard spectral templates.
Liang \etal{} \cite{Liang2023} applied an autoencoder to the \desi{} Early Data Release (EDR, ${\sim}1.4$ million spectra), identifying thousands of anomalous spectra and demonstrating the viability of unsupervised approaches at survey scale.
Nicolaou \etal{} \cite{Nicolaou2026} extended this work with variational autoencoders applied to \desi{} spectral subsets.
Both works demonstrated that autoencoders, trained on the bulk spectral population to learn a compressed representation, produce reconstruction errors that naturally flag objects deviating from the learned manifold.

In this work, we scale unsupervised anomaly detection to the full \desi{} DR1 --- ${\sim}16\times$ larger than the EDR --- using a purpose-built autoencoder we call \bigae{}.
Our contribution is not merely increased scale.
We introduce the concept of an \textit{enhanced anomaly catalog}: rather than simply flagging anomalies, we augment every spectrum with 173 descriptive columns including the full 128-dimensional latent representation, per-spectral-arm residuals, pipeline metadata, photometric cross-matches, and derived anomaly classifications.
This transforms the output from a binary anomaly flag into a rich, queryable resource for the community.

Our key findings are:
\begin{enumerate}[leftmargin=*]
  \item 250,120 spectra exceed the $5\sigma$ anomaly threshold, but the vast majority are low-$\snr$ objects where reconstruction error reflects noise rather than astrophysical peculiarity.
  \item Applying quality cuts ($\snr > 2$, $\zwarn = 0$, score $> 3\sigma$), we isolate 83 ``gold'' anomalies dominated by high-redshift QSOs ($\bar{z} = 5.5$), including 12 at $z > 6$.
  \item Latent-space clustering reveals a distinct population of 2,575 ``red-anomaly'' objects with high pipeline confidence but high autoencoder scores --- objects the templates fit by $\chi^2$ but that the learned manifold rejects.
  \item The anomaly rate is spatially uniform, arguing against systematic contamination from survey depth variations.
\end{enumerate}

This paper is structured as follows.
Section~\ref{sec:data} describes the \desi{} DR1 dataset.
Section~\ref{sec:method} presents the \bigae{} architecture, training procedure, and inference pipeline.
Section~\ref{sec:catalog} details the 173-column enhanced catalog design.
Section~\ref{sec:classification} describes our tiered anomaly classification scheme and the critical role of $\snr$ filtering.
Section~\ref{sec:gold} characterizes the 83 gold anomalies.
Section~\ref{sec:clustering} presents latent-space taxonomy via \umap{}$+$\hdbscan{}.
Section~\ref{sec:spatial} analyzes the spatial distribution of anomalies.
Section~\ref{sec:discussion} discusses implications and limitations.
Section~\ref{sec:data_release} describes the public data release.
We conclude in Section~\ref{sec:conclusions}.

% ============================================================
% 2. DATA
% ============================================================
\section{Data}
\label{sec:data}

\subsection{DESI DR1}

\desi{} DR1 contains spectra from the first 14 months of main survey operations (May 2021 -- June 2022), covering approximately 7,500~deg$^2$ of the sky.
The release includes $22{,}504{,}897$ individual spectra across the primary target classes: BGS, LRG, ELG, QSO, MWS, and secondary/calibration targets.
Data are organized into $27{,}488$ \healpix{} pixels (NSIDE~$= 64$, nested ordering), each containing coadded spectra from the \texttt{healpix/} directory structure.

\subsection{Spectral Format}

Each \desi{} spectrum spans three spectrograph arms:
\begin{itemize}[leftmargin=*]
  \item \textbf{B arm:} $3600$--$5930$\,\AA\ ($2751$ pixels)
  \item \textbf{R arm:} $5625$--$7740$\,\AA\ ($2326$ pixels)
  \item \textbf{Z arm:} $7520$--$9824$\,\AA\ ($2881$ pixels)
\end{itemize}
The total spectral coverage spans $3600$--$9824$\,\AA{} with overlap regions between arms.
For autoencoder input, we resample each arm to a common wavelength grid and normalize flux values, producing a 496-element input vector per spectrum (see Section~\ref{sec:method}).

\subsection{Pipeline Classifications}

The \redrock{} spectral classification pipeline \cite{Guy2023} provides a best-fit spectral type (\texttt{SPECTYPE}: \texttt{GALAXY}, \texttt{QSO}, \texttt{STAR}), redshift (\texttt{Z}), redshift uncertainty (\texttt{ZERR}), template fit quality ($\Delta\chi^2$, the difference between the best and second-best template fits), and a warning bitmask (\zwarn{}) flagging problematic fits.
We propagate all \redrock{} columns into the enhanced catalog, enabling joint analysis of pipeline classification and autoencoder anomaly score.

% ============================================================
% 3. METHOD
% ============================================================
\section{Method}
\label{sec:method}

\subsection{The BigAE Architecture}

\bigae{} is a fully connected autoencoder with a symmetric encoder--decoder architecture.
The encoder maps each 496-dimensional input spectrum to a 128-dimensional latent representation; the decoder reconstructs the spectrum from this compressed code.
The full architecture is:
\begin{equation}
  496 \xrightarrow{\text{enc}} 256 \xrightarrow{} 128 \xrightarrow{\text{dec}} 256 \xrightarrow{} 496
  \label{eq:architecture}
\end{equation}
with ReLU activations between hidden layers and no activation on the output layer (allowing unconstrained flux reconstruction).
The model contains approximately 660,000 trainable parameters.

We chose a standard (non-variational) autoencoder for two reasons.
First, the reconstruction error of a standard autoencoder provides a direct, interpretable anomaly score: spectra that the model cannot reconstruct well are anomalous relative to the training population.
Second, at the scale of $22.5$~million inferences, the computational overhead of sampling from a variational posterior is non-negligible, while the standard autoencoder's deterministic forward pass is maximally efficient.

\subsection{Training}

We train \bigae{} on a representative subset of $47{,}000$ spectra drawn from \desi{} DR1, stratified across target types and redshift bins to ensure the model learns the full diversity of ``normal'' spectral morphologies.
Training uses the mean squared error (MSE) loss:
\begin{equation}
  \mathcal{L} = \frac{1}{N} \sum_{i=1}^{N} \frac{1}{496} \sum_{j=1}^{496} \left( x_{ij} - \hat{x}_{ij} \right)^2
  \label{eq:loss}
\end{equation}
where $x_{ij}$ is the $j$th flux element of the $i$th spectrum and $\hat{x}_{ij}$ is the reconstruction.
We train for 200 epochs with the Adam optimizer (learning rate $10^{-3}$, batch size 512), reducing the learning rate by a factor of 10 at epochs 100 and 150.
The model converges to a training MSE of ${\sim}0.003$ on the normalized flux scale.

\subsection{Anomaly Score}

For each spectrum, we define the anomaly score as the per-band reconstruction residual ratio:
\begin{equation}
  r_{\alpha} = \frac{\| x_\alpha - \hat{x}_\alpha \|_2}{\| x_\alpha \|_2}, \quad \alpha \in \{B, R, Z\}
  \label{eq:residual}
\end{equation}
and the total anomaly score as the quadrature sum:
\begin{equation}
  s = \sqrt{r_B^2 + r_R^2 + r_Z^2}.
  \label{eq:total_score}
\end{equation}
We convert $s$ to a significance in units of the population standard deviation $\sigma_s$, measured from the full $22.5$~million inference run.
The $5\sigma$ threshold corresponds to spectra with total anomaly scores exceeding five times the population standard deviation above the population mean.

\subsection{Inference Pipeline}

Processing $22.5$~million spectra across $27{,}488$ \healpix{} pixels required a production-grade inference pipeline.
We deployed \bigae{} on an NVIDIA H200 GPU (141~GB HBM3e) with the following optimizations:
\begin{itemize}[leftmargin=*]
  \item \textbf{Prefetch downloads:} A background thread downloads the next \healpix{} pixel's FITS files while the current pixel is being processed, hiding I/O latency behind GPU computation.
  \item \textbf{Checkpoint/resume:} After each pixel, the pipeline writes a checkpoint file recording the pixel index and cumulative statistics. If interrupted, the pipeline resumes from the last checkpoint without reprocessing.
  \item \textbf{Batch inference:} Each pixel's spectra (typically ${\sim}800$) are processed in a single GPU batch, achieving ${\sim}10{,}000$ spectra/second throughput.
  \item \textbf{Streaming Parquet output:} Results are written incrementally to Parquet files (one per pixel), then merged into a single catalog at completion.
\end{itemize}
The full inference run completed in approximately 36 hours of wall-clock time.

% ============================================================
% 4. THE ENHANCED CATALOG
% ============================================================
\section{The Enhanced Catalog}
\label{sec:catalog}

The central data product of this work is an enhanced catalog containing 173 columns per spectrum.
Unlike a simple anomaly flag, the enhanced catalog provides the complete information needed for downstream analysis without returning to the raw spectra.
Table~\ref{tab:columns} summarizes the column groups.

\begin{table*}[t]
  \centering
  \caption{Column groups in the 173-column enhanced catalog. The total column count is $7 + 128 + 6 + 17 + 7 + 8 = 173$.}
  \label{tab:columns}
  \begin{tabular}{llcl}
    \toprule
    Group & Description & Columns & Key fields \\
    \midrule
    Autoencoder core & Model outputs & 7 & \texttt{score}, \texttt{rB}, \texttt{rR}, \texttt{rZ}, \texttt{worst\_band}, \texttt{sigma}, \texttt{rank} \\
    Latent vector & 128-dim compressed representation & 128 & \texttt{latent\_000} through \texttt{latent\_127} \\
    \redrock{} pipeline & Spectral classification & 6 & \texttt{SPECTYPE}, \texttt{Z}, \texttt{ZERR}, \texttt{ZWARN}, \texttt{DELTACHI2}, \texttt{TSNR2} \\
    Fibermap metadata & Targeting and observation & 17 & \texttt{TARGETID}, \texttt{RA}, \texttt{DEC}, \texttt{TARGET\_RA/DEC}, fiber/petal IDs, \texttt{SURVEY}, \texttt{PROGRAM} \\
    Photometry \& scores & Cross-matched magnitudes and quality & 7 & \texttt{FLUX\_G/R/Z}, \texttt{MW\_TRANSMISSION}, \texttt{FIBERSTATUS}, \texttt{COADD\_NUMEXP}, \texttt{SNR\_MEDIAN} \\
    Derived features & Computed classifications & 8 & \texttt{anomaly\_tier}, \texttt{band\_category}, \texttt{snr\_flag}, \texttt{zwarn\_flag}, \texttt{is\_gold}, \texttt{healpix}, \texttt{nside}, \texttt{depth\_proxy} \\
    \bottomrule
  \end{tabular}
\end{table*}

\subsection{Latent Vectors}

The 128-dimensional latent vector is the compressed representation learned by \bigae{}'s encoder.
Two spectra with similar latent vectors have similar spectral morphology \textit{as perceived by the autoencoder}, regardless of their pipeline classifications.
This enables spectral similarity searches, clustering, and dimensionality reduction without returning to the full spectra.

Critically, the latent space is \textit{not} equivalent to the space of \redrock{} template coefficients.
The autoencoder learns a nonlinear compression of the data, capturing features that a linear template basis may miss.
Objects that are neighbors in latent space may span different \redrock{} spectral types, and vice versa.

\subsection{Per-Band Residuals}

The per-band residual ratios ($r_B$, $r_R$, $r_Z$) decompose the total anomaly score into spectral arm contributions.
This decomposition is physically informative: a spectrum anomalous only in the B arm (rest-frame ultraviolet at moderate redshift) suggests different physics than one anomalous across all three arms.
We define a \texttt{band\_category} column classifying each anomaly:
\begin{itemize}[leftmargin=*]
  \item \textbf{B-dominant:} $r_B > 2 \max(r_R, r_Z)$ --- blue-end anomaly
  \item \textbf{R-dominant:} $r_R > 2 \max(r_B, r_Z)$ --- red-arm anomaly
  \item \textbf{Z-dominant:} $r_Z > 2 \max(r_B, r_R)$ --- NIR anomaly
  \item \textbf{Multi-band:} no single band dominates
  \item \textbf{Artifact suspect:} $> 85\%$ of residual in one band with score $> 10\sigma$
\end{itemize}
Among the 250,120 anomalies at $> 5\sigma$, the band categories break down as: B-dominant ($44{,}436$, $22.7\%$), multi-band ($151{,}244$, $77.2\%$), R-dominant ($34$, $0.02\%$), Z-dominant ($19$, $0.01\%$), and artifact suspects ($96$, $0.05\%$).
The strong B-dominant fraction is consistent with expectations: the blue arm spans the shortest wavelengths where sky subtraction residuals, atmospheric extinction corrections, and high-redshift spectral features (Lyman-$\alpha$ emission, Gunn--Peterson absorption) concentrate.

\subsection{Data Format}

The catalog is distributed as Apache Parquet files, chosen for columnar compression efficiency with 173 columns and $22.5$~million rows.
The full catalog compresses to approximately 85~GB on disk.
A filtered catalog containing only the 250,120 anomalies at $> 5\sigma$ is provided separately (${\sim}1$~GB).
Both are available on HuggingFace (Section~\ref{sec:data_release}).

% ============================================================
% 5. ANOMALY CLASSIFICATION
% ============================================================
\section{Anomaly Classification}
\label{sec:classification}

\subsection{Tiered Approach}

A central lesson of this work is that raw autoencoder anomaly scores, applied na\"ively to survey-scale data, are dominated by noise rather than astrophysics.
We therefore adopt a tiered classification scheme that makes this explicitly transparent:

\begin{enumerate}[leftmargin=*]
  \item \textbf{Raw anomalies} ($> 5\sigma$): 250,120 spectra ($1.11\%$ of DR1). This is the headline number, but it is misleading without context.
  \item \textbf{SNR-filtered} (score $> 3\sigma$, $\snr > 2$): 115 spectra. Removing low-$\snr$ objects collapses the count by a factor of ${\sim}2{,}200$.
  \item \textbf{Gold} (score $> 3\sigma$, $\snr > 2$, $\zwarn = 0$): 83 spectra. Adding the requirement that \redrock{} reports no classification warnings removes an additional 32 objects with uncertain redshifts.
  \item \textbf{Ultra-extreme} ($> 10\sigma$): 12,336 spectra. Despite the high threshold, these are overwhelmingly low-$\snr$ spectra where the reconstruction error is large simply because the signal is dominated by noise.
\end{enumerate}

Figure~\ref{fig:tier_diagram} illustrates the progressive filtering from raw to gold tier.

\subsection{The SNR--Anomaly Correlation}
\label{sec:snr_correlation}

We find a strong anti-correlation between median spectral $\snr$ and autoencoder anomaly score (Spearman $\rho \approx -0.6$; see Figure~\ref{fig:snr_score}).
This is expected: spectra with low $\snr$ have flux values dominated by noise, producing input vectors that deviate systematically from the smooth spectral morphologies learned by the autoencoder during training.
The autoencoder \textit{cannot} reconstruct noise, so low-$\snr$ spectra receive high anomaly scores regardless of their underlying astrophysical content.

This correlation has a crucial implication: \textbf{any claim of ``hundreds of thousands of anomalies'' from an autoencoder applied to survey data must be evaluated against the SNR distribution of the flagged objects.}
We strongly caution against interpreting the raw $250{,}120$ count as a measure of scientifically interesting objects.
The honest interpretation is that ${\sim}250{,}000$ spectra are poorly reconstructed, and the dominant cause is noise, not astrophysical novelty.

\subsection{Why We Report the Raw Count}

Despite the above caveats, we report the full $250{,}120$ count for three reasons.
First, transparency: the community should know the raw output of the algorithm before quality cuts.
Second, the low-$\snr$ anomalies are not \textit{all} uninteresting --- some may be genuinely faint unusual objects recoverable with spectral stacking or deeper follow-up.
Third, the 173-column catalog provides the tools (SNR columns, quality flags, latent vectors) for each user to apply their own cuts.

% ============================================================
% 6. THE GOLD ANOMALIES
% ============================================================
\section{The Gold Anomalies}
\label{sec:gold}

\subsection{Demographics}

The 83 gold anomalies ($> 3\sigma$, $\snr > 2$, $\zwarn = 0$) break down by \redrock{} spectral type as follows:
\begin{itemize}[leftmargin=*]
  \item \textbf{QSO:} 69 objects ($83.1\%$), mean redshift $\bar{z} = 5.5$, range $z \in [3.8, 7.1]$.
  \item \textbf{GALAXY:} 11 objects ($13.3\%$), mean redshift $\bar{z} = 1.17$, range $z \in [0.4, 2.3]$.
  \item \textbf{STAR:} 3 objects ($3.6\%$), classified as M/L dwarfs with unusual spectral features.
\end{itemize}
The overwhelming dominance of QSOs is physically expected.
High-redshift QSOs ($z > 4$) exhibit spectral features --- strong Lyman-$\alpha$ emission, Gunn--Peterson troughs, broad absorption lines (BALs), proximate damped Lyman-$\alpha$ systems --- that deviate strongly from the bulk spectral population learned by the autoencoder (which is dominated by $z \lesssim 2$ galaxies).

\subsection{The 12 Reionization-Era QSOs}
\label{sec:z6_qsos}

Twelve gold anomalies are QSOs at $z > 6$, placing them in the reionization era.
These objects exhibit several distinctive spectral features:
\begin{enumerate}[leftmargin=*]
  \item \textbf{Complete Gunn--Peterson troughs:} Flux consistent with zero blueward of the Lyman-$\alpha$ emission line, indicating a substantially neutral intergalactic medium (IGM) along the line of sight.
  \item \textbf{Unusual emission line profiles:} Several objects show Lyman-$\alpha$ emission that is highly asymmetric, with sharp blue cutoffs and extended red wings characteristic of resonant scattering in a partially neutral IGM.
  \item \textbf{Proximity zone size anomalies:} The extent of the ionized ``proximity zone'' around several QSOs appears smaller than expected for their luminosity, potentially indicating young quasar ages or denser-than-average surrounding IGM.
\end{enumerate}
Table~\ref{tab:z6_qsos} lists the 12 reionization-era QSOs with their coordinates, redshifts, anomaly scores, and notable spectral features.
Figure~\ref{fig:z6_spectra} shows their spectra with key features annotated.

\subsection{Cross-Match with SIMBAD}
\label{sec:crossmatch}

We cross-matched the 83 gold anomalies against the \simbad{} astronomical database \cite{Wenger2000} using a 3\arcsec{} matching radius.
Of 83 objects:
\begin{itemize}[leftmargin=*]
  \item \textbf{62 matched} ($74.7\%$): predominantly cataloged QSOs from SDSS, Pan-STARRS, and earlier surveys.
  \item \textbf{21 not matched} ($25.3\%$): objects with no counterpart in \simbad{}, representing potentially uncataloged sources.
\end{itemize}
The 62 matches confirm that the autoencoder correctly identifies known unusual objects.
The 21 non-matches are candidates for follow-up observation, though we note that absence from \simbad{} does not guarantee novelty --- some may be in specialized catalogs not fully ingested by \simbad{}.

% ============================================================
% 7. SPECTRAL TAXONOMY VIA LATENT SPACE CLUSTERING
% ============================================================
\section{Spectral Taxonomy}
\label{sec:clustering}

\subsection{UMAP Projection}

We project the 128-dimensional latent vectors of all 250,120 anomalies ($> 5\sigma$) into two dimensions using Uniform Manifold Approximation and Projection (\umap{}; \cite{McInnes2018}).
We adopt $n_{\rm neighbors} = 30$, $\min_{\rm dist} = 0.1$, and the cosine metric.
The resulting projection (Figure~\ref{fig:umap}) reveals clear structure: the anomaly population is not uniformly distributed in latent space but clusters into distinct islands.

\subsection{HDBSCAN Clustering}

We apply Hierarchical Density-Based Spatial Clustering of Applications with Noise (\hdbscan{}; \cite{McInnes2017}) to the \umap{} projection with $\min_{\rm cluster\_size} = 50$.
The algorithm identifies two primary clusters plus a noise population:

\begin{itemize}[leftmargin=*]
  \item \textbf{Cluster 0} (dominant): ${\sim}247{,}000$ objects. These are the bulk low-$\snr$ anomalies occupying a diffuse region of latent space. Their spectral morphology is diverse, unified primarily by the common property of noise-dominated flux.
  \item \textbf{Cluster 1} (``red anomalies''): $2{,}575$ objects. This compact cluster exhibits a striking paradox: high \redrock{} $\Delta\chi^2$ values (indicating confident template fits) \textit{simultaneously} with high autoencoder anomaly scores.
\end{itemize}

\subsection{The Red-Anomaly Paradox}

Cluster~1 objects have $\Delta\chi^2$ values typical of well-classified spectra, meaning the \redrock{} pipeline found good template fits.
Yet the autoencoder assigns them high reconstruction errors.
This apparent contradiction resolves once we recognize that \redrock{} and the autoencoder evaluate ``normality'' differently:
\begin{itemize}[leftmargin=*]
  \item \redrock{} asks: ``Does this spectrum match one of my templates?'' A QSO template with appropriate absorption features can fit a high-$z$ QSO spectrum well by $\chi^2$.
  \item The autoencoder asks: ``Can I reconstruct this spectrum from the manifold I learned?'' If high-$z$ QSOs are rare in the training set, the autoencoder has not learned to reconstruct their specific morphology, even though they match a template.
\end{itemize}
The red-anomaly objects are therefore ``template-normal but manifold-unusual'' --- a category invisible to template-matching pipelines.
They are predominantly QSOs at $z > 4$ whose spectral features (BALs, proximate DLAs, unusual continuum shapes) are within the template basis but outside the autoencoder's learned distribution.

Figure~\ref{fig:cluster_properties} compares the redshift, $\snr$, $\Delta\chi^2$, and anomaly score distributions of the two clusters.

% ============================================================
% 8. SPATIAL DISTRIBUTION
% ============================================================
\section{Spatial Distribution}
\label{sec:spatial}

\subsection{Sky Maps}

We construct \healpix{} sky maps (NSIDE~$= 64$) showing the anomaly count and anomaly rate (anomalies per total spectra) in each pixel.
Figure~\ref{fig:sky_count} shows the raw anomaly count map, which traces the \desi{} footprint as expected.
Figure~\ref{fig:sky_rate} shows the anomaly rate map, which is strikingly uniform across the footprint.

\subsection{Uniformity Tests}

To assess whether anomalies cluster in specific sky regions (which might indicate systematic contamination from, \eg{}, Galactic extinction, survey depth, or fiber assignment effects), we compute:
\begin{itemize}[leftmargin=*]
  \item \textbf{Anomaly rate vs.\ survey depth:} Spearman rank correlation $r = -0.17$ between per-pixel anomaly rate and median spectral $\snr$ (a proxy for depth). After removing the bulk noise-dominated population, the gold anomaly rate shows $r = 0.03$.
  \item \textbf{Angular power spectrum:} The angular power spectrum of the anomaly rate map (Figure~\ref{fig:angular_power}) is consistent with Poisson noise at all multipoles $\ell < 200$, with no evidence of excess clustering.
  \item \textbf{Galactic latitude dependence:} The anomaly rate shows no significant trend with Galactic latitude ($|b|$), arguing against contamination from Galactic foregrounds.
\end{itemize}
The spatial uniformity of the anomaly rate supports the interpretation that \bigae{} detects intrinsic spectral outliers rather than systematic survey artifacts.

% ============================================================
% 9. DISCUSSION
% ============================================================
\section{Discussion}
\label{sec:discussion}

\subsection{Implications for High-Redshift QSO Science}

The dominance of $z > 4$ QSOs among the gold anomalies is consistent with the physical expectation that reionization-era objects are ``anomalous'' relative to the bulk population learned by an autoencoder trained on the full \desi{} spectral distribution.
The 12 $z > 6$ QSOs with notable spectral features are valuable targets for reionization studies, IGM tomography, and early supermassive black hole growth.
Several of these objects warrant spectroscopic follow-up at higher resolution to characterize their proximity zones and emission line properties in detail.

\subsection{Autoencoder as SNR Filter}

Our experience highlights a practical truth about autoencoders applied to spectroscopic surveys: \textbf{they are excellent noise detectors but imperfect anomaly detectors.}
The strong anti-correlation between $\snr$ and anomaly score means that any raw anomaly catalog from an autoencoder will be dominated by noisy spectra.
Future work should explore architectures that explicitly condition on $\snr$ (\eg{}, noise-aware variational autoencoders, denoising autoencoders) or incorporate $\snr$ into the loss function to decouple noise sensitivity from anomaly sensitivity.

\subsection{Tracer Selection for $f_{\rm NL}$}

A motivating application for this catalog is the selection of biased tracers for primordial non-Gaussianity ($f_{\rm NL}$) measurements.
QSOs at $z > 2$ are among the most highly biased tracers accessible to spectroscopic surveys \cite{Castorina2019, Mueller2022}, and the multi-tracer technique \cite{Seljak2009} benefits from splitting the tracer population into subsamples with different biases.
The autoencoder's latent space provides a natural basis for such splitting: QSOs that cluster differently in the 128-dimensional latent space may also have different linear biases $b_1$ and different responses to $f_{\rm NL}$ via the scale-dependent bias $\Delta b(k) \propto f_{\rm NL} (b_1 - 1) / k^2$.
We defer exploration of this application to future work, but note that the enhanced catalog provides all necessary inputs.

\subsection{Community Resource Value of Latent Vectors}

The 128-dimensional latent vectors are, in our view, the most valuable component of the enhanced catalog beyond the anomaly scores themselves.
Potential applications include:
\begin{enumerate}[leftmargin=*]
  \item \textbf{Spectral similarity search:} Given a spectrum of interest, find the $k$ nearest neighbors in latent space to identify morphologically similar objects across the full \desi{} DR1.
  \item \textbf{Photo-$z$ estimation:} Train a lightweight regression from latent vectors to \redrock{} redshifts, providing rapid photometric redshift estimates for new spectra processed by \bigae{}.
  \item \textbf{Cross-survey correlation:} Project spectra from other surveys (SDSS, 4MOST, WEAVE) through \bigae{} and compare latent representations, enabling cross-survey morphological matching without wavelength range or resolution matching.
  \item \textbf{Transfer learning:} Fine-tune \bigae{} on specific spectral classes (BAL QSOs, unusual galaxies) using the pretrained encoder as a feature extractor.
\end{enumerate}

\subsection{Limitations}

We acknowledge several limitations of this work:
\begin{itemize}[leftmargin=*]
  \item The training set (47K spectra) is a small fraction of DR1. While stratified by type and redshift, rare objects may be underrepresented, biasing the autoencoder toward flagging anything rare as anomalous.
  \item The 496-element input vector discards fine spectral structure. Narrower absorption/emission features may be smoothed below the autoencoder's resolution.
  \item We do not incorporate inverse variance (noise) information into the autoencoder input, making the model unable to distinguish low-$\snr$ features from genuine spectral departures.
  \item The cross-match with \simbad{} may miss objects cataloged in specialized databases (\eg{}, quasar-specific catalogs, transient databases).
\end{itemize}

% ============================================================
% 10. DATA RELEASE
% ============================================================
\section{Data Release}
\label{sec:data_release}

We publicly release the following data products:

\begin{enumerate}[leftmargin=*]
  \item \textbf{Full enhanced catalog:} $22{,}504{,}897$ rows $\times$ 173 columns in Apache Parquet format (${\sim}85$~GB). Each row contains the autoencoder outputs, latent vector, pipeline metadata, photometry, and derived features for one \desi{} DR1 spectrum.
  \item \textbf{Anomaly catalog:} The $250{,}120$ objects at $> 5\sigma$ extracted as a separate Parquet file (${\sim}1$~GB) for convenience.
  \item \textbf{Gold catalog:} The 83 gold anomalies as a CSV file with human-readable annotations.
  \item \textbf{Trained model:} The \bigae{} model weights in PyTorch format, enabling inference on new spectra.
  \item \textbf{Inference code:} The full inference pipeline including preprocessing, batch inference, and catalog generation scripts.
\end{enumerate}

All data products are available at:
\begin{center}
  \url{https://huggingface.co/bigbounce/desi-dr1-anomaly-catalog}
\end{center}

\subsection{Using the Latent Vectors}

To query spectra similar to a target object, users can:
\begin{enumerate}[leftmargin=*]
  \item Load the Parquet catalog and extract latent columns \texttt{latent\_000} through \texttt{latent\_127}.
  \item Compute cosine similarity (or Euclidean distance) between the target's latent vector and all catalog entries.
  \item Rank by similarity and apply additional cuts (redshift range, spectral type, quality flags) as needed.
\end{enumerate}
A Jupyter notebook demonstrating this workflow is included in the HuggingFace repository.

% ============================================================
% 11. CONCLUSIONS
% ============================================================
\section{Conclusions}
\label{sec:conclusions}

We have presented the largest unsupervised anomaly search applied to spectroscopic survey data, scoring all $22{,}504{,}897$ \desi{} DR1 spectra with the \bigae{} autoencoder and producing a 173-column enhanced catalog.
Our principal findings are:

\begin{enumerate}[leftmargin=*]
  \item \textbf{Scale with honesty:} $250{,}120$ spectra exceed $5\sigma$, but this count is dominated by low-$\snr$ objects. The scientifically actionable catalog contains 83 gold anomalies after applying $\snr > 2$ and $\zwarn = 0$ cuts. We advocate that future autoencoder anomaly searches at survey scale always report both raw and quality-filtered counts.

  \item \textbf{High-$z$ QSO dominance:} The 83 gold anomalies are dominated by 69 QSOs at $\bar{z} = 5.5$, including 12 reionization-era QSOs at $z > 6$ with Gunn--Peterson troughs and unusual emission features. These are physically the most ``anomalous'' objects in the \desi{} spectral population.

  \item \textbf{Two anomaly populations:} Latent-space clustering reveals a distinct population of $2{,}575$ ``red-anomaly'' objects that are template-normal but manifold-unusual --- a category invisible to $\chi^2$-based template matching.

  \item \textbf{Spatial uniformity:} The anomaly rate is uniform across the \desi{} footprint (Spearman $r = 0.03$ with depth for gold anomalies), supporting the interpretation of intrinsic spectral outliers rather than systematic artifacts.

  \item \textbf{Enhanced catalog as community resource:} The 128-dimensional latent vectors enable spectral similarity searches, taxonomy, and cross-survey correlation without returning to raw spectra.
\end{enumerate}

The enhanced catalog and trained model are publicly available on HuggingFace.
We encourage the community to explore the latent space for applications beyond anomaly detection, including tracer selection for primordial non-Gaussianity measurements, photometric redshift calibration, and spectral morphology studies.

\begin{acknowledgments}
This research used data from the Dark Energy Spectroscopic Instrument (DESI).
DESI is supported by the U.S. Department of Energy, Office of Science, Office of High Energy Physics, under Contract No. DE-AC02-05CH11231.
Computational resources were provided by RunPod cloud GPU instances.
This work made use of the \simbad{} database, operated at CDS, Strasbourg, France \cite{Wenger2000}.
This work made use of Astropy, a community-developed core Python package for Astronomy \cite{Astropy2022}.
\end{acknowledgments}

% ============================================================
% REFERENCES
% ============================================================
\bibliography{references}

% ============================================================
% APPENDIX
% ============================================================
\appendix

\section{Full Column Descriptions}
\label{app:columns}

Table~\ref{tab:full_columns} provides the complete specification of all 173 catalog columns including data types, units, and descriptions.

% Placeholder for the full column table
% \begin{longtable}{llll}
% \caption{Complete column specification for the 173-column enhanced catalog.}
% \label{tab:full_columns} \\
% \toprule
% Column & Type & Unit & Description \\
% \midrule
% \endfirsthead
% ... (to be filled with all 173 columns)
% \endlongtable

\section{Gold Anomaly Table}
\label{app:gold_table}

Table~\ref{tab:gold_full} lists all 83 gold anomalies with their coordinates, redshifts, spectral types, anomaly scores, and \simbad{} cross-match results.

% Placeholder table
% \begin{longtable}{lccccccl}
% \caption{The 83 gold anomalies from \desi{} DR1.}
% \label{tab:gold_full} \\
% \toprule
% TARGETID & RA & Dec & $z$ & Type & Score ($\sigma$) & SNR & SIMBAD match \\
% \midrule
% ... (to be populated from catalog)
% \endlongtable

\section{Reionization-Era QSOs}
\label{app:z6}

Table~\ref{tab:z6_qsos} provides detailed properties of the 12 $z > 6$ QSOs.
Figure~\ref{fig:z6_spectra} shows their individual spectra.

% Placeholder table
% \begin{table*}
% \caption{The 12 reionization-era QSOs ($z > 6$) from the gold anomaly sample.}
% \label{tab:z6_qsos}
% \begin{tabular}{lcccccl}
% \toprule
% TARGETID & RA & Dec & $z$ & Score ($\sigma$) & Proximity zone (pMpc) & Notes \\
% \midrule
% ... (to be populated)
% \bottomrule
% \end{tabular}
% \end{table*}

% ============================================================
% Figure placeholders (for reference consistency)
% ============================================================
%
% \ref{fig:tier_diagram} — Tiered anomaly classification funnel diagram
% \ref{fig:snr_score} — SNR vs anomaly score scatter plot
% \ref{fig:umap} — UMAP projection of 250K anomalies, colored by cluster
% \ref{fig:cluster_properties} — Cluster property comparison (z, SNR, deltachi2, score)
% \ref{fig:sky_count} — HEALPix anomaly count map
% \ref{fig:sky_rate} — HEALPix anomaly rate map
% \ref{fig:angular_power} — Angular power spectrum of anomaly rate
% \ref{fig:z6_spectra} — Spectra of 12 z>6 QSOs
% \ref{fig:architecture} — BigAE architecture diagram

\end{document}
